% ============================================================================
% Seccion 2 — Descripcion del dataset (punto 2 del enunciado)
% ============================================================================
\section{Descripci\'on del dataset}

\subsection{Ficha t\'ecnica}

\begin{center}
\begin{tabular}{p{4.2cm}p{10.5cm}}
\toprule
\textbf{Nombre} & Denuncias por desaparici\'on de personas \\
\textbf{Organizaci\'on} & Ministerio del Interior del Per\'u (MININTER) \\
\textbf{Rubro PNDA} & Seguridad y Justicia \\
\textbf{Fuente} & Plataforma Nacional de Datos Abiertos
  (\url{https://www.datosabiertos.gob.pe/dataset/denuncias-por-desaparici\%C3\%B3n-de-personas-mininter}) \\
\textbf{Recurso analizado} & \texttt{BD\_DESAPARICIONES\_ENERO\_2019\_DICIEMBRE\_2024.csv} (8.1 MB) \\
\textbf{Actualizaci\'on del recurso} & 18 de setiembre de 2026 \\
\textbf{Cobertura} & enero 2019 -- diciembre 2024 (72 meses) \\
\textbf{Grano} & a\~no $\times$ mes $\times$ ubigeo del hecho $\times$ perfil de la v\'ictima \\
\textbf{Registros} & 93,831 filas \\
\textbf{Columnas} & 13 \\
\textbf{Medida} & \texttt{CANTIDAD}: personas denunciadas en ese cruce (1 a 18; total 121,427) \\
\textbf{Formato} & CSV separado por comas, codificaci\'on UTF-8 con BOM, fin de l\'inea CRLF \\
\textbf{Recursos complementarios} & diccionario de datos (\texttt{.xlsx}) y metadatos (\texttt{.docx}) \\
\bottomrule
\end{tabular}
\end{center}

\subsection{Teor\'ia que implican los datos}
El registro no contiene personas identificables: contiene \emph{hechos de desaparici\'on
agregados}. Cada fila es el cruce de una fecha (a\~no y mes), una ubicaci\'on del hecho (ubigeo,
departamento, provincia, distrito) y un perfil de la v\'ictima (sexo, rango de edad, estado
civil, nacionalidad, nivel educativo y ocupaci\'on), con la cantidad de personas denunciadas en
ese cruce como medida. Es, por tanto, un cubo de datos en forma plana: el hecho es la
desaparici\'on, las dimensiones describen el cu\'ando, el d\'onde y el qui\'en, y la medida es
aditiva, lo que permite totalizar por cualquiera de las dimensiones sin recalcular el dataset.

Desde la pol\'itica p\'ublica, los datos responden a la teor\'ia de la \emph{b\'usqueda urgente}:
la probabilidad de encontrar a una persona desaparecida cae r\'apidamente con las horas, por lo
que el inter\'es anal\'itico se concentra en la distribuci\'on territorial (d\'onde se denuncia
m\'as), en el perfil etario (las y los adolescentes son el grupo m\'as expuesto) y en la
estacionalidad (meses con m\'as carga operativa para las comisar\'ias).

\subsection{Campos, cardinalidad y valores vac\'ios}

\begin{center}
\small
\begin{tabular}{llrp{6.6cm}}
\toprule
\textbf{Campo} & \textbf{Tipo} & \textbf{Valores} & \textbf{Descripci\'on y observaciones}\\
\midrule
A\~NO & entero & 6 & 2019 a 2024. Distribuci\'on de registros: 18,934 / 14,373 / 15,728 / 14,916 / 14,607 / 15,273.\\
MES & entero & 12 & 1 a 12. Los 72 cruces a\~no-mes tienen registros.\\
UBIGEO\_HECHO & texto & 1,506 & C\'odigo de 6 d\'igitos (departamento-provincia-distrito), con ceros a la izquierda que se conservan como texto. Un mismo ubigeo (150121) aparece con dos nombres de distrito: PUEBLO LIBRE y MAGDALENA VIEJA (su nombre antiguo).\\
DPTO\_HECHO & texto & 26 & Incluye ``LIMA METROPOLITANA'' y ``REGI\'ON LIMA'' como unidades separadas, y ``CALLAO'' como departamento.\\
PROV\_HECHO & texto & 195 & Nombre de la provincia del hecho.\\
DIST\_HECHO & texto & 1,506 & Nombre del distrito del hecho.\\
SEXO & texto & 3 & F (54,595), M (39,162) y 74 vac\'ios (0.08\,\%).\\
RANGO\_EDAD & texto & 7 & 0-5, 6-11, 12-17 (45,303 registros, el 48.3\,\%), 18-29, 30-59, 60 a m\'as y 1,243 vac\'ios (1.3\,\%).\\
EST\_CIVIL & texto & 7 & Soltero(a) 63,583; NO INDICA 8,286; casado(a) 6,147; conviviente, divorciado(a), viudo(a); 13,480 vac\'ios (14.4\,\%).\\
NACIONALIDAD & texto & 57 & Per\'u 91,171 (97.2\,\%), Venezuela 2,118, Colombia 123, Ecuador 70, Chile 54.\\
NIVEL\_EDUCATIVO & texto & 6 & Secundaria 21,782, primaria 9,195, superior 1,157, t\'ecnica, sin instrucci\'on; 60,731 vac\'ios (64.7\,\%).\\
OCUPACION & texto & 289 & Estudiante 20,894 (22.3\,\%), ama de casa 2,808, independiente 1,272; 60,790 vac\'ios (64.8\,\%).\\
CANTIDAD & entero & 18 & Personas denunciadas en el registro: 1 (78,128 registros, 83.3\,\%), 2, 3, \dots hasta 18 como m\'aximo.\\
\bottomrule
\end{tabular}
\end{center}

\subsection{Reglas de los datos identificadas}
\begin{enumerate}
  \item \textbf{La medida es positiva y acotada.} \texttt{CANTIDAD} toma valores entre 1 y 18; no
        hay ceros ni negativos, por lo que la suma de la columna es directamente el total de
        personas (121,427).
  \item \textbf{Un registro no equivale a una persona.} El 83.3\,\% de las filas tiene
        \texttt{CANTIDAD} = 1; contar filas (93,831) subestimar\'ia el fen\'omeno en 27,596
        personas. Toda consulta del informe agrega con \texttt{SUM(CANTIDAD)}.
  \item \textbf{El ubigeo es un identificador textual y tiene una excepci\'on.} Tiene ceros a la
        izquierda (por ejemplo \texttt{021809} para Nuevo Chimbote, \'Ancash); leerlo como
        n\'umero lo corromper\'ia. El ubigeo \texttt{150121} aparece con dos nombres de distrito
        (PUEBLO LIBRE y su antiguo nombre, MAGDALENA VIEJA), lo que obliga a colapsar la
        dimensi\'on geogr\'afica a una fila por ubigeo antes de hacer joins.
  \item \textbf{Las ausencias son informativas, no aleatorias.} Nivel educativo y ocupaci\'on
        faltan en cerca del 65\,\% de los registros, y faltan juntos (mismo patr\'on de 60,731 y
        60,790 vac\'ios), lo que sugiere un formulario que solo se completa en algunas
        comisar\'ias. No se imput\'o: se etiquet\'o como \texttt{NO ESPECIFICA}.
  \item \textbf{Nacionalidad no es sin\'onimo de migrante.} ``PER\'U'' es un valor m\'as de la
        columna; para analizar poblaci\'on extranjera hay que excluirlo expl\'citamente.
  \item \textbf{El territorio es jer\'arquico.} Los 1,506 ubigeos se agrupan en 195 provincias y
        26 unidades departamentales, lo que permite clasificarlas por regi\'on natural (costa,
        sierra, selva) sin informaci\'on externa.
  \item \textbf{No hay duplicados exactos de la llave completa.} Las 93,831 filas son
        combinaciones distintas de fecha, ubigeo y perfil.
\end{enumerate}

\subsection{Identidad del dataset y no duplicidad}
El registro se publica en la PNDA como un \emph{dataset independiente}: tiene su propia URL,
su propio identificador en el portal, su diccionario de datos (\texttt{.xlsx}) y su documento de
metadatos (\texttt{.docx}), y no figura como subconjunto declarado de otro dataset del cat\'alogo.
Su sistema de origen es el SIDPOL del MININTER, que alimenta tambi\'en otros registros policiales
publicados por separado (denuncias policiales, indicadores de seguridad ciudadana); lo que se
analiza aqu\'i es el registro de desapariciones publicado como tal, con su propia cobertura
(enero 2019 -- diciembre 2024) y sus 13 campos.

Se deja constancia de dos decisiones para que el trabajo sea trazable:
\begin{itemize}
  \item Se usa \textbf{un solo recurso} (el de 13 columnas, 2019--2024). El portal publica
        adem\'as un recurso 2019--2025 de 10 columnas que, en los seis a\~nos comparables, suma
        exactamente lo mismo por a\~no (0.00 de diferencia). Apilar ambos duplicar\'ia todos los
        casos, de modo que se eligi\'o el recurso con mayor detalle de columnas.
  \item No se imputan valores ni se excluyen filas por tener campos vac\'ios; esas ausencias se
        documentan como regla del dato (secci\'on 2.4) y se conservan como categor\'ia
        \texttt{NO ESPECIFICA}.
\end{itemize}

\subsection{Informaci\'on que se puede obtener}
Con este registro se pueden responder, entre otras, las siguientes l\'ineas de an\'alisis, que
son las que estructuran las 25 consultas del laboratorio:
\begin{itemize}
  \item \textbf{Concentraci\'on territorial}: qu\'e distritos, provincias y departamentos acumulan
        m\'as personas desaparecidas y cu\'al es su peso relativo.
  \item \textbf{Perfil de la v\'ictima}: distribuci\'on por sexo, grupo etario, estado civil,
        nivel educativo y ocupaci\'on, y combinaciones entre ellos.
  \item \textbf{Vulnerabilidad de menores}: proporci\'on de menores de 18 a\~nos por departamento
        y comparaci\'on con la poblaci\'on adulta.
  \item \textbf{Tendencia y estacionalidad}: evoluci\'on anual, variaci\'on interanual mes a mes e
        \'indice de cada mes frente al esperado.
  \item \textbf{Migraci\'on}: peso y evoluci\'on de las personas de nacionalidad extranjera.
  \item \textbf{Calidad y cobertura}: meses o territorios sin registros, integridad referencial
        del cat\'alogo geogr\'afico y efecto del periodo pand\'emico (2020--2021).
\end{itemize}
